%==============================================================================
% SHABNAM HAKIMI — CURRICULUM VITAE
% Built with Awesome-CV (https://github.com/posquit0/Awesome-CV)
%==============================================================================

\documentclass[11pt, a4paper]{awesome-cv}

%------------------------------------------------------------------------------
% PACKAGES
%------------------------------------------------------------------------------
\usepackage{etoolbox}

%------------------------------------------------------------------------------
% BIBLIOGRAPHY STYLE
%
% Change the style name below. Common options:
%   numeric      — [1], [2], … (default)
%   ieee         — numbered, abbreviated authors & journal names
%   nature       — numbered superscripts, condensed
%   apa          — author-year (disable showBibNumbers / reverseBibNumbers)
%   authoryear   — generic author-year
%   chicago-authordate
%   mla
%
% NOTE: showBibNumbers and reverseBibNumbers only work with numeric-family
% styles (numeric, ieee, nature). For author-year styles, also set:
%   \togglefalse{showBibNumbers}
%------------------------------------------------------------------------------
\PassOptionsToPackage{style=numeric}{biblatex}   % ← change style here

\usepackage[backend=biber,
            sorting=none]{biblatex}
\usepackage{xstring}
\usepackage[symbol,perpage]{footmisc}  % symbol footnotes, reset per page
\usepackage{perpage}                   % gives us \MakePerPage for custom resets

%------------------------------------------------------------------------------
% PUBLICATION ANNOTATION SYMBOLS
%
% \cvfnmark{equalcontrib} — outputs a superscript * to mark equal-contribution
% authors in usera fields. Always outputs *, no dynamic counter needed.
%
% The legend ("* equal contributions...") is printed once at the top of the
% Publications section via \cvpublegend, not repeated per subsection.
%------------------------------------------------------------------------------
\newcommand{\cvfnmark}[1]{\textsuperscript{*}}

% \cvpublegend: prints the symbol legend once at the top of a section.
% Place immediately after \cvsection{Publications}.
\newcommand{\cvpublegend}{%
  \vspace{-2pt}%
  {\footnotesize * equal contributions, shared first authorship}%
  \vspace{4pt}%
}

%==============================================================================
% BIBLIOGRAPHY NUMBERING TOGGLES & MACHINERY
%==============================================================================
%
% \toggletrue{showBibNumbers}    → show [1], [2], … labels
% \togglefalse{showBibNumbers}   → hide labels entirely
%
% \toggletrue{reverseBibNumbers} → N…1 (most recent = highest number)
% \togglefalse{reverseBibNumbers}→ 1…N (forward order)
%
% Numbering restarts at 1 (or N) for each bibliography section.
%
% Implementation: we use our own per-item counter entirely independent of
% biblatex's internal numbering. \AtEveryBibitem increments a running index;
% we write each section's total to the .aux so the next pass can compute
% reverse labels. No refsection environments needed.
%==============================================================================
\newtoggle{showBibNumbers}
\newtoggle{reverseBibNumbers}

\toggletrue{showBibNumbers}     % true = show numbers | false = hide numbers
\togglefalse{reverseBibNumbers} % true = reverse (N…1) | false = forward (1…N)

% Per-section counters
\newcounter{cvbibindex}   % increments with each printed entry (reset per section)
\newcounter{cvbibtotal}   % total entries in current section (loaded from aux)

% cvbibindex starts at 1. \AtEveryBibitem fires AFTER the label is rendered,
% so the label reads the current value (1, 2, 3…) and then the increment
% advances it for the next entry.
\AtEveryBibitem{\stepcounter{cvbibindex}}

% Override the printed label to use our own counter.
% labelnumber controls the number inside brackets in the entry body.
% We also override the list \item label via the labelnumberwidth format
% and the blx@lbx numeric label macro.
\DeclareFieldFormat{labelnumber}{%
  \iftoggle{showBibNumbers}{%
    \iftoggle{reverseBibNumbers}{%
      \number\numexpr\value{cvbibtotal}-\value{cvbibindex}+1\relax%
    }{%
      \number\value{cvbibindex}%
    }%
  }{%
    \hspace{-\labelsep}%
  }%
}

% Also patch the label used for \item[...] in the bibliography list env.
% In biblatex's numeric style this is \printtext[labelnumberwidth]{...}
% which calls \mkbibbrackets{\printfield{labelnumber}} — so overriding
% labelnumber above is sufficient.

% When numbers are hidden, remove the label gutter entirely
\AtBeginDocument{%
  \nottoggle{showBibNumbers}{%
    \setlength{\biblabelsep}{0pt}%
    \defbibenvironment{bibliography}
      {\list{}
         {\setlength{\leftmargin}{\bibhang}%
          \setlength{\itemindent}{-\leftmargin}%
          \setlength{\itemsep}{\bibitemsep}%
          \setlength{\parsep}{\bibparsep}}}
      {\endlist}
      {\item}%
  }{}%
}

% \startbibsection{name} — call before each \printbibliography block.
%   Resets the index counter and loads the stored total for this section.
%   On the first compile the total is 0 (labels may be wrong); correct on
%   subsequent compiles once \finishbibsection has written the total to .aux.
\newcommand{\startbibsection}[1]{%
  \setcounter{cvbibindex}{1}%
  \ifcsdef{cvbibtotal@#1}{%
    \setcounter{cvbibtotal}{\csname cvbibtotal@#1\endcsname}%
  }{%
    \setcounter{cvbibtotal}{0}%
  }%
}

% \finishbibsection stores (cvbibindex - 1) as the total, since cvbibindex
% ends up one past the last entry after the final increment.
\newwrite\cvbibcountsfile
\AtBeginDocument{\immediate\openout\cvbibcountsfile=cvbibcounts.tex}
\AtEndDocument{\immediate\closeout\cvbibcountsfile}

\newcommand{\finishbibsection}[1]{%
  \immediate\write\cvbibcountsfile{%
    \string\expandafter\string\def
    \string\csname\space cvbibtotal@#1\string\endcsname
    {\number\numexpr\value{cvbibindex}-1\relax}%
  }%
}

% Load counts written by previous pass
\AtBeginDocument{\IfFileExists{cvbibcounts.tex}{\expandafter\def\csname cvbibtotal@journals\endcsname{16}
\expandafter\def\csname cvbibtotal@preprints\endcsname{7}
\expandafter\def\csname cvbibtotal@conference\endcsname{14}
\expandafter\def\csname cvbibtotal@presentations\endcsname{54}
\expandafter\def\csname cvbibtotal@bookchapters\endcsname{0}
\expandafter\def\csname cvbibtotal@scicomm\endcsname{1}
\expandafter\def\csname cvbibtotal@patents\endcsname{11}
}{}}

%------------------------------------------------------------------------------
% BOLD OWN NAME IN BIBLIOGRAPHY
%
% Entries that go through the normal name formatter (no usera field)
% will have "Hakimi" bolded automatically. Entries with a usera field
% bypass the formatter — bold is applied manually via \textbf in the .bib.
% Change \cvownfamilyname if the CV owner's name changes.
%------------------------------------------------------------------------------
\def\cvownfamilyname{Hakimi}

\DeclareNameFormat{boldifown}{%
  \ifstrequal{\namepartfamily}{\cvownfamilyname}{%
    \cvbibname{%
      \usebibmacro{name:family-given}{\namepartfamily}{\namepartgiveni}%
                                     {\namepartprefix}{\namepartsuffix}%
    }%
  }{%
    \usebibmacro{name:family-given}{\namepartfamily}{\namepartgiveni}%
                                   {\namepartprefix}{\namepartsuffix}%
  }%
  \usebibmacro{name:andothers}%
}

% NOTE: entries using the `usera` field bypass the name formatter.
% Bold your name manually there, e.g.: authorlist = {\textbf{Hakimi, S}, ...}

%------------------------------------------------------------------------------
% CUSTOM AUTHOR STRING — uses built-in `usera` field
%
% For entries with equal-contribution authors or other non-standard author
% formatting, set the `usera` field in the .bib file to a pre-formatted
% LaTeX string. Biber always passes `usera` through as a literal (no parsing).
%
% Use \cvbibname{} (not \textbf{}) for bold names — it uses \bfseries which
% works reliably in Awesome-CV's fontspec context. Use ~ for spaces inside
% name initials, e.g. "Hakimi,~S" to prevent TeX from dropping the space.
%
% Example:
%   usera = {\cvbibname{Hakimi,~S}\cvfnmark{equalcontrib}, Lee,~J\cvfnmark{equalcontrib}},
%
% When `usera` is present it replaces the normal \printnames{author} output.
% When absent, \printnames{author} runs normally (with boldifown applied).
%------------------------------------------------------------------------------

% \cvbibname: bold name for usera fields and boldifown.
% Keeps \sffamily (matching bib body font) but forces series b (Bold),
% overriding the \fontseries{l} (Light) set by \bodyfontlight.
\newcommand{\cvbibname}[1]{{\sffamily\fontseries{b}\selectfont #1}}

\DeclareFieldFormat{usera}{#1}

\renewbibmacro*{author}{%
  \iffieldundef{usera}{%
    \printnames{author}%
  }{%
    \printfield{usera}%
  }%
}

% Apply boldifown alias at begin-document so it overrides biblatex's own
% default aliases (which are set during \begin{document}).
\AtBeginDocument{%
  \DeclareNameAlias{author}{boldifown}%
  \DeclareNameAlias{editor}{boldifown}%
}

% Note: the note field prints normally for all entry types.
% Equal-contribution markers should be in the authorlist field (via \cvfnmark),
% NOT in the note field. See journals.bib and conference.bib for examples.
% Entries with note = {* equal contributions...} should have that note removed.

\addbibresource{refs/journals.bib}
\addbibresource{refs/preprints.bib}
\addbibresource{refs/conference.bib}
\addbibresource{refs/patents.bib}
\addbibresource{refs/scicomm.bib}
\addbibresource{refs/presentations.bib}

\nocite{*}  % include all entries from all .bib files

%------------------------------------------------------------------------------
% BIBLATEX HEADING — cvsubsection style
% Renders bibliography section titles using \cvsubsection (Awesome-CV).
% Usage: heading=cvsubsection  (replaces heading=subbibliography)
%------------------------------------------------------------------------------
\defbibheading{cvsubsection}[\bibname]{%
  \cvsubsection{#1}%
}

%------------------------------------------------------------------------------
% BIBLATEX CUSTOM CHECKS — patent status × selected toggle
%------------------------------------------------------------------------------

% Exclude entries tagged keyword=presentation from preprints section
\defbibcheck{notpresentation}{%
  \ifkeyword{presentation}{\skipentry}{}%
}
\defbibcheck{grantedpatents}{%
  \ifentrytype{patent}{%
    \iffieldequalstr{status}{Granted}{}{\skipentry}%
  }{\skipentry}%
}

\defbibcheck{grantedpatentsselected}{%
  \ifentrytype{patent}{%
    \iffieldequalstr{status}{Granted}{%
      \ifkeyword{selected}{}{\skipentry}%
    }{\skipentry}%
  }{\skipentry}%
}

\defbibcheck{filedpatents}{%
  \ifentrytype{patent}{%
    \iffieldequalstr{status}{Filed}{}{\skipentry}%
  }{\skipentry}%
}

\defbibcheck{filedpatentsselected}{%
  \ifentrytype{patent}{%
    \iffieldequalstr{status}{Filed}{%
      \ifkeyword{selected}{}{\skipentry}%
    }{\skipentry}%
  }{\skipentry}%
}

%------------------------------------------------------------------------------
% PATENT PRINT COMMAND
%------------------------------------------------------------------------------
\newcommand{\printpatents}{%
  \ifthenelse{\equal{\patentMode}{none}}{%
    % Print nothing
  }{%
    \ifthenelse{\equal{\patentMode}{granted}}{%
      \ifthenelse{\equal{\patentFilter}{selected}}{%
        \leavevmode\printbibliography[check=grantedpatentsselected,
                           title={Granted},
                           heading=cvsubsection]%
      }{%
        \leavevmode\printbibliography[check=grantedpatents,
                           title={Granted},
                           heading=cvsubsection]%
      }%
    }{%
      % All (granted + filed)
      \ifthenelse{\equal{\patentFilter}{selected}}{%
        \leavevmode\printbibliography[check=grantedpatentsselected,
                           title={Granted},
                           heading=cvsubsection]%
        \leavevmode\printbibliography[check=filedpatentsselected,
                           title={Filed},
                           heading=cvsubsection]%
      }{%
        \leavevmode\printbibliography[check=grantedpatents,
                           title={Granted},
                           heading=cvsubsection]%
        \leavevmode\printbibliography[check=filedpatents,
                           title={Filed},
                           heading=cvsubsection]%
      }%
    }%
  }%
}

%==============================================================================
% TOGGLES — SECTION VISIBILITY
%==============================================================================
\newtoggle{showJournals}
\newtoggle{showPreprints}
\newtoggle{showConferenceProceedings}
\newtoggle{showPatents}
\newtoggle{showSciComm}
\newtoggle{showBookChapters}
\newtoggle{showTalks}
\newtoggle{showConferencePresentations}
\newtoggle{showTeaching}
\newtoggle{showService}
\newtoggle{showReviewer}
\newtoggle{showGrants}
\newtoggle{showMentorship}
\newtoggle{showMemberships}
\newtoggle{showReferences}
\newtoggle{showFullReferences}

\toggletrue{showJournals}
\toggletrue{showPreprints}
\toggletrue{showConferenceProceedings}
\toggletrue{showPatents}
\toggletrue{showSciComm}
\toggletrue{showBookChapters}
\toggletrue{showTalks}
\toggletrue{showConferencePresentations}
\toggletrue{showTeaching}
\toggletrue{showService}
\toggletrue{showReviewer}
\toggletrue{showGrants}
\toggletrue{showMentorship}
\toggletrue{showMemberships}
\toggletrue{showReferences}
\togglefalse{showFullReferences}  % false = "Available upon request"

%==============================================================================
% TOGGLES — PUBLICATION FILTERS
%==============================================================================
\newtoggle{selectedOnlyPublications}    % true = "Selected Publications" section title
\newtoggle{selectedOnlyPatents}         % true = "Selected Patents" section title
\newtoggle{selectedOnlyJournals}
\newtoggle{selectedOnlyConferenceProceedings}
\newtoggle{selectedOnlyPreprints}
\newtoggle{selectedOnlyPresentations}

\togglefalse{selectedOnlyPublications}  % true = "Selected Publications" heading
\togglefalse{selectedOnlyPatents}       % true = "Selected Patents" heading (mirrors patentFilter)
\togglefalse{selectedOnlyJournals}
\togglefalse{selectedOnlyConferenceProceedings}
\togglefalse{selectedOnlyPreprints}
\toggletrue{selectedOnlyPresentations}

%------------------------------------------------------------------------------
% PATENT DISPLAY MODE
%
% \def\patentMode{none}       → omit patents section entirely
% \def\patentMode{granted}    → show granted patents only
% \def\patentMode{all}        → show granted + filed patents
%
% \def\patentFilter{selected} → only entries with keywords={selected}
% \def\patentFilter{all}      → all entries regardless of keyword
%------------------------------------------------------------------------------
\def\patentMode{granted}     % none | granted | all
\def\patentFilter{selected}  % selected | all

%==============================================================================
% TOGGLES — SECTION DEPTH
%==============================================================================
\newtoggle{showAllGrants}
\newtoggle{showAllTalks}
\newtoggle{showAllTeaching}
\newtoggle{showAllService}

\togglefalse{showAllGrants}
\togglefalse{showAllTalks}
\togglefalse{showAllTeaching}
\togglefalse{showAllService}

% First-pass guard: skip all \printbibliography calls until the .bbl exists.
% This prevents xelatex from crashing before biber has run, which would
% produce a malformed .bcf that biber then refuses to process.
\newif\ifbblexists
\IfFileExists{\jobname.bbl}{\bblexiststrue}{\bblexistsfalse}

% Wrap \printbibliography to be a no-op on the first pass
\let\origprintbibliography\printbibliography
\renewcommand{\printbibliography}[1][]{%
  \ifbblexists\origprintbibliography[#1]\fi%
}

% Note: \leavevmode before each \printbibliography ensures LaTeX is in
% horizontal mode, preventing \spacefactor crashes from biblatex warnings.

%------------------------------------------------------------------------------
% AWESOME-CV CONFIGURATION
%------------------------------------------------------------------------------
\geometry{left=1.4cm, top=.8cm, right=1.4cm, bottom=1.8cm, footskip=.5cm}

\colorlet{awesome}{black}
\colorlet{darktext}{black}
\setbool{acvSectionColorHighlight}{false}

%------------------------------------------------------------------------------
% PERSONAL INFO
%------------------------------------------------------------------------------
\name{Shabnam}{Hakimi, Ph.D.}
\position{Senior Research Scientist{\enskip\cdotp\enskip}Human-Centered AI}
\address{Toyota Research Institute, 4440 El Camino Real, Los Altos, CA 94022}
\email{shabnamhakimi@gmail.com}
\homepage{www.shabnamhakimi.com}
% \github{yourgithub}
% \linkedin{yourlinkedin}
% \orcid{0000-0000-0000-0000}

%------------------------------------------------------------------------------
% DOCUMENT
%------------------------------------------------------------------------------
\begin{document}

\makecvheader[C]

\makecvfooter
  {\today}
  {Shabnam Hakimi ~~·~~ Curriculum Vitae}
  {\thepage}

%------------------------------------------------------------------------------
% SECTIONS
%------------------------------------------------------------------------------

%==============================================================================
% EDUCATION
%==============================================================================
\cvsection{Education}

\begin{cventries}

  \cventry
    {Ph.D., Computation and Neural Systems}
    {California Institute of Technology}
    {Pasadena, CA}
    {2014}
    {}

  \cventry
    {B.A., Psychology}
    {Stanford University}
    {Stanford, CA}
    {2006}
    {}

\end{cventries}
%==============================================================================
% PROFESSIONAL EXPERIENCE
%==============================================================================
\cvsection{Professional Experience}
\begin{cventries}
  \cventry
    {Senior Research Scientist}
    {Toyota Research Institute}
    {Los Altos, CA}
    {01/2021 -- Present}
    {\,}
  \cventry
    {Postdoctoral Scholar}
    {Center for Cognitive Neuroscience, Duke University}
    {Durham, NC}
    {06/2016 -- 01/2021}
    {\,}
  \cventry
    {Postdoctoral Research Associate}
    {Institute of Cognitive Science, University of Colorado}
    {Boulder, CO}
    {04/2015 -- 05/2016}
    {\,}
  \cventry
    {Director}
    {CaféWell Behavioral Labs, Welltok}
    {Denver, CO}
    {08/2013 -- 01/2015}
    {\,}
  \cventry
    {Graduate Research and Teaching Fellow}
    {Computation and Neural Systems, California Institute of Technology}
    {Pasadena, CA}
    {09/2008 -- 08/2013}
    {\,}
  \cventry
    {Intramural Research Training Award Fellow}
    {Genes, Cognition \& Psychosis Program, National Institute of Mental Health (NIH)}
    {Bethesda, MD}
    {08/2006 -- 07/2008}
    {\,}
  \cventry
    {Research Assistant}
    {Department of Psychology, Stanford University}
    {Stanford, CA}
    {05/2003 -- 06/2006}
    {\,}
\end{cventries}

\iftoggle{showGrants}{%==============================================================================
% GRANTS & AWARDS
%
% TOGGLE GUIDE (set in main.tex):
%   \toggletrue{showAllGrants}  → show ALL entries (full CV)
%   \togglefalse{showAllGrants} → show only selected entries
%==============================================================================

\iftoggle{showAllGrants}
  {\cvsection{Grants \& Awards}}
  {\cvsection{Selected Grants \& Awards}}

\begin{cventries}

  % ── SELECTED (always shown) ────────────────────────────────────────────────

  \cventry
    {Ruth Kirchstein National Research Service Award (NRSA)}
    {National Institute of Mental Health, NIH}
    {Bethesda, MD}
    {2017 -- 2020}
    {}
    % {\begin{cvitems}
    %   \item {\textit{Project}: ``Modeling the neurobehavioral basis of extrinsic and volitional motivation in ADHD''}
    %   \item {\textit{Role}: PI (F32 MH114608)}
    % \end{cvitems}}

  \cventry
    {Scientific Research Network on Decision Neuroscience \& Aging (SRNDNA) Pilot Grant}
    {National Institute of Aging, NIH}
    {Bethesda, MD}
    {2019 -- 2020}
    {}
    % {\begin{cvitems}
    %   \item {\textit{Project}: ``Neurobehavioral correlates of volitional regulation of motivation across the lifespan''}
    %   \item {\textit{Role}: Subaward PI (R24 AG054355 to G.R. Samanez-Larkin)}
    % \end{cvitems}}

  \cventry
    {SciComm Fellow, Duke Initiative for Science and Society}
    {Duke University}
    {Durham, NC}
    {2016 -- 2017}
    {}

  \cventry
    {NSF Integrative Graduate Education and Research Traineeship (IGERT) Fellow}
    {California Institute of Technology}
    {Pasadena, CA}
    {2008 -- 2013}
    {}

  \cventry
    {Intramural Research Training Award}
    {National Institute of Mental Health, NIH}
    {Bethesda, MD}
    {2006 -- 2008}
    {}

  % ── FULL CV ONLY ──────────────────────────────────────────────────────────
  \iftoggle{showAllGrants}{%
    \cventry
      {Neurohackweek (Neurohackacademy) Fellow}
      {University of Washington eScience Institute}
      {Seattle, WA}
      {2016}
      {}

    \cventry
      {Application of Science Award Runner-Up}
      {NeuroLeadership Institute}
      {New York, NY}
      {2016}
      {}

    \cventry
      {Summer Institute in Cognitive Neuroscience Fellow}
      {University of California, Santa Barbara}
      {Santa Barbara, CA}
      {2009}
      {}

    \cventry
      {NSF/Caltech IGERT fMRI Award}
      {California Institute of Technology}
      {Pasadena, CA}
      {2009}
      {}

    \cventry
      {HealthEmotions Research Institute Symposium Scholar}
      {University of Wisconsin Symposium on Emotion}
      {Madison, WI}
      {2008}
      {}

    \cventry
      {HealthEmotions Research Institute Symposium Fellow}
      {University of Wisconsin Symposium on Emotion}
      {Madison, WI}
      {2007}
      {}

    \cventry
      {HealthEmotions Research Institute Symposium Scholar}
      {University of Wisconsin Symposium on Emotion}
      {Madison, WI}
      {2006}
      {}

    \cventry
      {Undergraduate Research Programs Conference Travel Grant}
      {Stanford University}
      {Stanford, CA}
      {2006}
      {}

    \cventry
      {Vice Provost for Undergraduate Education Summer Research Grant}
      {Stanford University}
      {Stanford, CA}
      {2005}
      {}

    \cventry
      {HealthEmotions Research Institute Symposium Scholar}
      {University of Wisconsin Symposium on Emotion}
      {Madison, WI}
      {2005}
      {}

    \cventry
      {Vice Provost for Undergraduate Education Summer Research Grant}
      {Stanford University}
      {Stanford, CA}
      {2004}
      {}

    % Add additional grants/awards below this line:
  }{}

\end{cventries}}{}

%-- Publications --------------------------------------------------------------
% Top-level heading: "Publications" or "Selected Publications"
% Subsections: Journals · Preprints · Conference Proceedings ·
%              Conference Presentations · Book Chapters
%
% Toggle \toggletrue{selectedOnlyPublications} to switch the section title.
% Individual subsection selected-only filters remain independent.

\newbool{showPublications}
\boolfalse{showPublications}
\iftoggle{showJournals}{\booltrue{showPublications}}{}
\iftoggle{showPreprints}{\booltrue{showPublications}}{}
\iftoggle{showConferenceProceedings}{\booltrue{showPublications}}{}
\iftoggle{showConferencePresentations}{\booltrue{showPublications}}{}
\iftoggle{showBookChapters}{\booltrue{showPublications}}{}

\ifbool{showPublications}{%
  \iftoggle{selectedOnlyPublications}{%
    \cvsection{Selected Publications}%
  }{%
    \cvsection{Publications}%
  }%
  \cvpublegend%

  \iftoggle{showJournals}{%
    \startbibsection{journals}%
    \iftoggle{selectedOnlyJournals}{%
      \leavevmode\printbibliography[type=article, keyword=selected,
                         title={Peer-Reviewed Journal Publications (Selected)},
                         heading=cvsubsection]%
    }{%
      \leavevmode\printbibliography[type=article,
                         title={Peer-Reviewed Journal Publications},
                         heading=cvsubsection]%
    }%
    \finishbibsection{journals}%
  }{}

  \iftoggle{showPreprints}{%
    \startbibsection{preprints}%
    \iftoggle{selectedOnlyPreprints}{%
      \leavevmode\printbibliography[type=unpublished, keyword=selected,
                         check=notpresentation,
                         title={Preprints \& Working Papers (Selected)},
                         heading=cvsubsection]%
    }{%
      \leavevmode\printbibliography[type=unpublished,
                         check=notpresentation,
                         title={Preprints \& Working Papers},
                         heading=cvsubsection]%
    }%
    \finishbibsection{preprints}%
  }{}

  \iftoggle{showConferenceProceedings}{%
    \startbibsection{conference}%
    \iftoggle{selectedOnlyConferenceProceedings}{%
      \leavevmode\printbibliography[type=inproceedings, keyword=selected,
                         title={Peer-Reviewed Conference Proceedings (Selected)},
                         heading=cvsubsection]%
    }{%
      \leavevmode\printbibliography[type=inproceedings,
                         title={Peer-Reviewed Conference Proceedings},
                         heading=cvsubsection]%
    }%
    \finishbibsection{conference}%
  }{}

  \iftoggle{showConferencePresentations}{%
    \startbibsection{presentations}%
    \iftoggle{selectedOnlyPresentations}{%
      \leavevmode\printbibliography[type=unpublished, keyword=presentation, keyword=selected,
                         title={Selected Conference Presentations},
                         heading=cvsubsection]%
    }{%
      \leavevmode\printbibliography[type=unpublished, keyword=presentation,
                         title={Conference Presentations},
                         heading=cvsubsection]%
    }%
    \finishbibsection{presentations}%
  }{}

  \iftoggle{showBookChapters}{%
    \startbibsection{bookchapters}%
    \leavevmode\printbibliography[type=incollection,
                       title={Reviews, Book Chapters \& Commentaries},
                       heading=cvsubsection]%
    \finishbibsection{bookchapters}%
  }{}
}{}

%-- Patents -------------------------------------------------------------------

\iftoggle{showPatents}{%
  \ifthenelse{\equal{\patentMode}{none}}{}{%
    \iftoggle{selectedOnlyPatents}{%
      \cvsection{Selected Patents}%
    }{%
      \cvsection{Patents}%
    }%
    \startbibsection{patents}%
    \vspace{2pt}%
    \printpatents
    \finishbibsection{patents}%
  }%
}{}

%-- Talks & Science Communication ---------------------------------------------

\newbool{showTalksSection}
\boolfalse{showTalksSection}
\iftoggle{showTalks}{\booltrue{showTalksSection}}{}
\iftoggle{showSciComm}{\booltrue{showTalksSection}}{}

\ifbool{showTalksSection}{%
  \cvsection{Talks \& Science Communication}%
  \vspace{2pt}%
  % sections/talks.tex should use \cvsubsection{Invited Talks} (not \cvsection)
  % to render as a subsection under this top-level heading.
  \iftoggle{showTalks}{\begin{cventries}

  \cventry
    {``Psychologically-guided generative AI interventions for flexible, creative decision making.''}
    {West Coast Neuroeconomics Mini-Symposium, University of California, Los Angeles}
    {Los Angeles, CA}
    {2025}
    {}

    {``Abstract semantic prompts shape sequential decision making in design.''}
    {West Coast Neuroeconomics Mini-Symposium, University of California, San Diego}
    {San Diego, CA}
    {2025}
    {}

  \cventry
    {``Neuroforecasting market demand for electrified vehicles: Leveraging the neuroscience of affect and motivation to predict future consumer behaviors.''}
    {2025 Annual Convention of the Western Psychological Association}
    {Las Vegas, NV}
    {2025}
    {}

  \cventry
    % {\textbf{Hakimi, S}, Hong, M, \& Klenk, M. ``Resolving mental simulation of possible futures for understanding and supporting design.''}
    {``Resolving mental simulation of possible futures for understanding and supporting design.''}
    {Design Computing \& Cognition '24 Workshop: ``Rethinking How `Design Spaces' are Constructed''}
    {Montréal, Canada}
    {2024}
    {}

  \cventry
    {``Rethinking user understanding and applications for innovation in design.''}
    {Berkeley Institute of Design Seminar, University of California, Berkeley}
    {Berkeley, CA}
    {2024}
    {}

  \cventry
    {``Learning to self-regulate motivation using real-time fMRI neurofeedback.''}
    {Department of Psychology, Temple University}
    {Philadelphia, PA}
    {2018}
    {}

  \cventry
    {``Learning to regulate motivation using real-time neurofeedback.''}
    {Department of Economics, University of Zürich}
    {Zürich, Switzerland}
    {2018}
    {}

  \iftoggle{showAllTalks}{

    \cventry
      {``The academic, there and back again.''}
      {Center for Cognitive Neuroscience Colloquium, Duke University}
      {Durham, NC}
      {2019}
      {}

    \cventry
      {``Modeling the neural basis of individual differences in intertemporal choice.''}
      {Central Institute of Mental Health (ZI)}
      {Mannheim, Germany}
      {2016}
      {}

    \cventry
      {``Neural correlates of intertemporal choice.''}
      {Department of Computer Science, University of Colorado Boulder}
      {Boulder, CO}
      {2015}
      {}

  }{}

\end{cventries}}{}

  \iftoggle{showSciComm}{%
    \startbibsection{scicomm}%
    \leavevmode\printbibliography[type=misc, keyword=scicomm,
                       title={Science Communication},
                       heading=cvsubsection]%
    \finishbibsection{scicomm}%
  }{}%
}{}

%-- Teaching & Service --------------------------------------------------------
\iftoggle{showTeaching}{%==============================================================================
% TEACHING
%==============================================================================

\iftoggle{showAllTeaching}
  {\cvsection{Teaching}}
  {\cvsection{Selected Teaching}}

\begin{cventries}

  \cventry
    {Guest Lecturer}
    {Duke University}
    {Durham, NC}
    {Fall 2018}
    {\begin{cvitems}
      \item {\textit{``Motivation and Decision Making''}}
      \item {DECSCI 101 / PSY 141: Fundamentals of Decision Science}
    \end{cvitems}}

  \cventry
    {Guest Lecturer \& Curriculum Designer}
    {California Institute of Technology}
    {Pasadena, CA}
    {Fall 2011}
    {\begin{cvitems}
      \item {CNS/Bi/Ph 107: Writing About Scientific Research}
    \end{cvitems}}

  \cventry
    {Teaching Assistant}
    {California Institute of Technology}
    {Pasadena, CA}
    {Fall 2009}
    {\begin{cvitems}
      \item {Bi/CNS 150: Introduction to Neurobiology}
    \end{cvitems}}

  % ── FULL CV ONLY ──────────────────────────────────────────────────────────
  \iftoggle{showAllTeaching}{%
    \cventry
      {Teaching Assistant}
      {Wellcome Trust Conference Center}
      {Hinxton, Cambridge, UK}
      {Summer 2009}
      {\begin{cvitems}
        \item {Wellcome Trust School on the Biology of Social Cognition}
      \end{cvitems}}%
    % Add additional teaching entries below this line:
  }{}

\end{cventries}}{}
\iftoggle{showService}{%==============================================================================
% PROFESSIONAL SERVICE & OUTREACH
%==============================================================================

\newcommand{\cventryselected}[5]{\cventry{#1}{#2}{#3}{#4}{#5}}

\cvsection{Professional Service \& Outreach}

\begin{cventries}

  % ── 2022 ──────────────────────────────────────────────────────────────────

  \cventryselected
    {Co-Chair, Women \& Allies Employee Resource Group (WAERG) \textnormal{\textit{(elected)}}}
    {Toyota Research Institute}
    {Los Altos, CA}
    {03/2022 -- 03/2024}
    {}

  % ── 2021 ──────────────────────────────────────────────────────────────────

  \iftoggle{showAllService}{
    \cventry
      {Idea Review Mentor \& Pitch Coach}
      {Project Invent}
      {}
      {10/2021 -- 04/2022}
      {}
    \vspace{-13.5pt}
  }{}

  \cventryselected
    {Postdoctoral Representative, Center for Cognitive Neuroscience \textnormal{\textit{(nominated \& elected)}}}
    {Duke University}
    {Durham, NC}
    {08/2019 -- 01/2021}
    {}

  \cventryselected
    {Co-Founder \& Leadership Committee Member}
    {Duke Inclusion and Power Dynamics in Academia Series}
    {Durham, NC}
    {11/2017 -- 01/2021}
    {\begin{cvitems}
      \item {Co-I: Duke Faculty Advancement Seed Grants (2018--2020)}
    \end{cvitems}}

  \cventryselected
    {Co-Founder \& Leadership Committee Member}
    {SPEAK: Scientists Promoting Equity and Knowledge}
    {Durham, NC}
    {06/2016 -- 01/2021}
    {}

  \cventryselected
    {Journal Club Panel Member, \textit{PNAS} Front Matter}
    {National Academy of Sciences}
    {Washington, DC}
    {06/2018 -- 08/2020}
    {}

  \cventryselected
    {Communications Committee Member \& Student/Postdoc SIG Liaison}
    {Organization for Human Brain Mapping}
    {Minneapolis, MN}
    {01/2017 -- 12/2018}
    {}

  \cventryselected
    {Guest Instructor \& Curriculum Designer, Course: \textit{Decision Making}}
    {Carolina Friends School}
    {Durham, NC}
    {03/2017 -- 05/2017}
    {}

  % ── 2012 ──────────────────────────────────────────────────────────────────

  \iftoggle{showAllService}{
    \cventry
      {Presenter, Reel Science at Caltech}
      {California Institute of Technology}
      {Pasadena, CA}
      {03/2012 -- 04/2013}
      {}
    \vspace{-13.5pt}
  }{}

  % ── 2006 ──────────────────────────────────────────────────────────────────

  \iftoggle{showAllService}{
    \cventry
      {Clinical Volunteer, Behavioral Health \& Ortho/Neuro Trauma Units}
      {Suburban Hospital}
      {Bethesda, MD}
      {10/2006 -- 07/2007}
      {}

    % Add additional service entries below this line:

  }{}

\end{cventries}}{}
\iftoggle{showReviewer}{%==============================================================================
% REVIEWING
%==============================================================================
\cvsection{Reviewing}

\cvsubsection{Journals}\par
\begin{cvparagraph}
\textit{Cognitive, Affective, and Behavioral Neuroscience}; \textit{Computational Psychiatry}; \textit{Experimental Brain Research}; \textit{Journal of Child Psychiatry and Psychology}; \textit{Journal of Experimental Psychology: Learning, Memory, and Cognition}; \textit{Memory and Cognition}; \textit{Nature Human Behaviour}; \textit{NeuroImage}; \textit{PLoS One}; \textit{Psychological Science}; \textit{Social, Cognitive, and Affective Neuroscience}
\end{cvparagraph}

\cvsubsection{Conferences}\par
\begin{cvparagraph}
ACM CHI Conference on Human Factors in Computing Systems (2026); ACM Creativity and Cognition Conference (C\&C, 2026); Cognitive Science Society (2025--2026); Generative AI and HCI Workshop at CHI (GenAICHI, 2025); Organization for Human Brain Mapping (OHBM; 2010--2011, 2016--2019); Society for Neuroeconomics (SNE, 2025--2026); Triangle Society for Neuroscience (SfN, 2018); Workshop on Behavioral ML at NeurIPS (2024)
\end{cvparagraph}}{}

\end{document}